\chapter{Introduction}

\section{Background and Motivation}
Medical education is undergoing a structural transformation driven by competency-based curricula, rising enrollment, and limited bedside exposure. While high-stakes licensing examinations such as the USMLE continue to rely on large-scale multiple-choice assessment, clinical competence additionally requires history-taking, physical examination reasoning, diagnostic synthesis, documentation (SOAP notes), and communication under time pressure---skills traditionally acquired through clerkships and standardized patient (OSCE) encounters. These training modalities are expensive, scheduling-intensive, and difficult to scale across institutions.

Concurrent advances in large language models (LLMs) have enabled conversational agents that simulate patient personas, explain clinical reasoning, and generate structured educational feedback at scale \cite{topol_highperf_2019,wartman_med_ed_2019}. Commercial question banks (e.g., UWorld-style interfaces) excel at drill-based knowledge retention but offer limited authentic encounter practice. Institutions, meanwhile, require dashboards to assign cases, monitor completion, compare cohort performance, and align training with curricular objectives. A fragmented toolchain---separate LMS, messaging, billing, analytics, and AI tutors---increases integration cost and data silos.

\textbf{MedPrepAI} addresses this gap as a unified \emph{clinical-lab} platform: an intelligent virtual consultation, medical training, and clinical workflow system that spans the entire repository scope---not only the four MedPrep modes, but also question banks, mock exams, subscriptions, faculty tooling, chat, notifications, study groups, achievements, DOCX ingestion, deployment infrastructure, and administrative governance. The project title reflects this breadth: virtual consultation (AI patient interviews), medical training (modes, QBank, flashcards, AI tutor), and clinical workflow (assignments, institution linking, faculty analytics, payments, audit trails).

The proliferation of AI in healthcare education also raises requirements for \emph{governed} access: students should receive mode-appropriate scaffolding (Learning vs.\ independent Practice), evaluation should be rubric-aligned, and institutions must control entitlements through subscriptions and roles. MedPrepAI embeds these policies in backend guards and frontend route protection, ensuring that AI capabilities remain pedagogically purposeful rather than unconstrained chat.

\section{Problem Statement}
Despite availability of QBank products and isolated virtual patient demos, there is no single open, extensible platform that simultaneously provides:
\begin{enumerate}
  \item Multi-mode AI clinical encounters with persistent transcripts and OSCE-style documentation;
  \item Institution-grade workflow (assignments, faculty analytics, domain-linked cohorts);
  \item Monetizable feature gating aligned with medical school licensing models;
  \item Production engineering (modular API, migrations, realtime messaging, payment webhooks).
\end{enumerate}
MedPrepAI is engineered to close this integration gap within one codebase (\texttt{clinical-lab}), enabling capstone-level demonstration of both AI innovation and software systems design.

\section{Goals and Objectives}
The primary goal of this project is to design, implement, and evaluate \textbf{MedPrepAI} as a production-grade platform that unifies AI-augmented clinical encounters with LMS, institutional workflow, and commercial subscription infrastructure.

The specific objectives are:
\begin{enumerate}
  \item To provide \textbf{safe, scalable virtual consultation} through four MedPrep modes (Practice, Learning, Evaluation, Shadow), enabling students to interview AI patients, order workups, submit diagnoses, and complete SOAP notes with automated scoring and hint telemetry.
  \item To deliver a \textbf{comprehensive medical training stack} including hierarchical content (products, sections, chapters, topics), rich question authoring (Markdown and DOCX pipelines), assessments, mock exams, flashcards, bookmarks, notes, study plans, goals, and an AI Tutor grounded in medical pedagogy.
  \item To support \textbf{clinical and institutional workflow}: faculty workspaces, student comparison, assignments, institution-domain linking, discussions, study groups, feedback tickets, and question reporting---aligned with how medical colleges operate cohort-based training.
  \item To implement \textbf{secure, monetizable access control} combining JWT authentication, RBAC, permission matrices, subscription packages, Stripe payments, and per-feature entitlements (including per-mode MedPrep access).
  \item To architect a \textbf{maintainable full-stack system} using NestJS, Prisma (multi-file MySQL schema), Next.js, WebSockets for realtime features, Swagger-documented APIs, and cloud deployment---demonstrating software engineering practices expected of a BS CIS capstone.
  \item To leverage \textbf{LLMs (Google Gemini)} for patient responses, faculty shadow reasoning, differential updates, case generation, SOAP grading, and tutoring---with structured prompts, persistence in \texttt{MedprepConversation} records, and fallback paths when API keys are unavailable.
\end{enumerate}

\section{Functional and Non-Functional Requirements}
Table~\ref{tab:requirements} summarizes requirements derived from stakeholder needs (students, faculty, administrators).

\begin{table}[htbp]
  \centering
  \caption{Functional requirements of MedPrepAI.}
  \label{tab:requirements}
  \small
  \begin{tabular}{@{}p{1.2cm}p{4.2cm}p{7.8cm}@{}}
    \toprule
    \textbf{ID} & \textbf{Category} & \textbf{Requirement} \\
    \midrule
    FR-01 & Virtual patient & System shall simulate patient dialogue constrained by case metadata. \\
    FR-02 & Modes & System shall support Practice, Learning, Evaluation, and Shadow with distinct UX. \\
    FR-03 & Persistence & All encounter messages, SOAP, and diagnoses shall be stored in MySQL. \\
    FR-04 & QBank & Users with entitlement shall browse, filter, and test on tagged questions. \\
    FR-05 & Authoring & Admins shall create/import questions via UI, Markdown, or DOCX. \\
    FR-06 & Faculty & Faculty shall assign cases and view student MedPrep metrics. \\
    FR-07 & Billing & Students shall purchase subscriptions; features unlock via entitlements. \\
    FR-08 & Messaging & Users shall participate in chat rooms and receive notifications. \\
    FR-09 & Tutor & Subscribed users shall converse with AI medical tutor. \\
    FR-10 & Analytics & Platform shall track achievements, goals, and study-plan progress. \\
    \bottomrule
  \end{tabular}
\end{table}

\begin{table}[htbp]
  \centering
  \caption{Non-functional requirements.}
  \label{tab:nfr}
  \small
  \begin{tabular}{@{}p{1.2cm}p{3.5cm}p{8.5cm}@{}}
    \toprule
    \textbf{ID} & \textbf{Attribute} & \textbf{Requirement} \\
    \midrule
    NFR-01 & Security & Passwords hashed; JWT validated; logout blacklists tokens. \\
    NFR-02 & Scalability & Stateless API; indexed queries on user/session hot paths. \\
    NFR-03 & Maintainability & Modular Prisma schemas; Nest feature modules; typed DTOs. \\
    NFR-04 & Observability & Audit logs; structured API errors; demo seed for QA. \\
    NFR-05 & Deployability & Documented build/rsync/PM2/Nginx procedures. \\
    NFR-06 & Safety & LLM prompts disclaim real-patient advice; educational use only. \\
    \bottomrule
  \end{tabular}
\end{table}

\section{Contributions}
Table~\ref{tab:contributions} lists principal contributions of this thesis relative to a minimal student project.

\begin{table}[htbp]
  \centering
  \caption{Principal contributions of the MedPrepAI project.}
  \label{tab:contributions}
  \begin{tabular}{@{}p{3.8cm}p{9.2cm}@{}}
    \toprule
    \textbf{Contribution} & \textbf{Description} \\
    \midrule
    Unified clinical-lab platform & Single monorepo integrating AI encounters, LMS, workflow, and commerce. \\
    Four-mode pedagogy & Distinct training modalities with shared persistence schema and entitlement slugs. \\
    Entitlement engine & Parameterized \texttt{EntitlementDefinition} + \texttt{valueJson} beyond boolean features. \\
    Shadow-mode observability & Replay, interventions, structured lab reports synced to SHADOW sessions. \\
    Faculty operations layer & Compare students, assignments, institution sync from MedPrep completion. \\
    Content pipelines & Markdown/DOCX ingestion for scalable QBank population. \\
    Engineering rigor & 25+ API modules, guards, migrations, deployment guide, demo seed covering all modes. \\
    \bottomrule
  \end{tabular}
\end{table}

\section{Scope of the Clinical-Lab Repository}
This thesis documents the complete \texttt{clinical-lab} monorepo, including:
\begin{itemize}
  \item \textbf{backend/}: NestJS API modules (auth, users, roles, products, content, questions, assessments, subscriptions, payments, medprep-ai, ai-tutor, faculty, institution, chat, notifications, achievements, discussions, study-groups, goals, feedback, mock-exams, flashcards, bookmarks, notes, study-plans, student-stats, realtime).
  \item \textbf{frontend-next/}: Next.js 15 application with landing page, dashboards, MedPrep UIs, faculty pages, payments, test creation, org-chart, db-viewer, and API routes proxying LLM calls.
  \item \textbf{backend/prisma/schema/}: Modular database definitions spanning 20 schema files.
  \item Supporting assets: deployment guides, access-control documentation, DOCX parser workflows, demo seed scripts, and legacy UI references (\texttt{uworld-replit}, \texttt{uworld-bolt}).
\end{itemize}

Out of scope: regulatory certification as a medical device, integration with live hospital EHR production systems, and multilingual clinical validation studies.

\section{Stakeholders}
Table~\ref{tab:stakeholders-intro} identifies primary stakeholders; expanded analysis appears in Appendix~\ref{app:stakeholders}.

\begin{table}[htbp]
  \centering
  \caption{Primary stakeholders of MedPrepAI.}
  \label{tab:stakeholders-intro}
  \begin{tabular}{@{}lp{10cm}@{}}
    \toprule
    \textbf{Stakeholder} & \textbf{Interest} \\
    \midrule
    Medical students & Safe AI practice, exam preparation, progress tracking \\
    Clinical faculty & Assignments, analytics, messaging \\
    Institution leadership & Scalable licensing, cohort outcomes \\
    System administrators & User management, content pipelines, revenue \\
    \bottomrule
  \end{tabular}
\end{table}

\section{Development Methodology}
The project followed an iterative agile workflow:
\begin{itemize}
  \item \textbf{Phase 1 --- Foundation}: Auth, roles, Prisma multi-schema, product/content models (LMS skeleton).
  \item \textbf{Phase 2 --- QBank UX}: Port UWorld-style UI; assessments, test sessions, explanations.
  \item \textbf{Phase 3 --- MedPrep AI}: Practice/Learning flows, Gemini integration, session persistence.
  \item \textbf{Phase 4 --- Shadow \& Evaluation}: Attending replay, rubric scoring, hint telemetry.
  \item \textbf{Phase 5 --- Institution \& commerce}: Faculty module, Stripe, entitlements, assignments.
  \item \textbf{Phase 6 --- Community \& polish}: Discussions, study groups, achievements, landing page, deployment.
\end{itemize}

\section{Thesis Organization}
Chapter~\ref{ch:lit} reviews literature. Chapter~\ref{ch:method} presents architecture and implementation (including diagrams and API tables). Chapter~\ref{ch:results} reports case studies. Chapter~\ref{ch:conclusion} concludes. Appendices list modules, routes, and configuration.
